% practice.tex -- Module 8 (From the Universe to the Atom): Practice set.
% 25 new HSC-style questions: Section I 9 multiple choice, Section II 10 short answer,
% Section III 6 extended response. 81 marks. Ordered from easy to HSC-hard.
% Independent check (solved again without the answers):
% Answers: practice_answers.tex (one practiceans per practiceq, same order).
% Local helpers (prefix "prac"), defined with \providecommand so the answers file can
% define them too: \pracmarks, \pracblank, \pracNuc, \pracaq.
% Figures and tables here use \caption* (no number): this chapter is unnumbered, so a
% numbered caption would continue the numbering of Chapter 5.
\providecommand\pracmarks[1]{\hfill\mbox{\textcolor{hscmuted}{\small(#1)}}}
% \pracblank: an empty cell for the student to fill in (a short grey line).
\providecommand\pracblank{\textcolor{hscmuted}{\rule[-2pt]{2.2em}{0.4pt}}}
% \pracNuc{A}{Z}{X}: nuclide symbol, A above Z, both right-aligned before X.
\providecommand\pracNuc[3]{\ensuremath{\mathord{%
  \setbox0=\hbox{$\scriptstyle#1$}\setbox2=\hbox{$\scriptstyle#2$}%
  \ifdim\wd0<\wd2 \setbox0=\hbox to\wd2{\hss\box0}\else\setbox2=\hbox to\wd0{\hss\box2}\fi
  {}^{\box0}_{\box2}\mathrm{#3}}}}
% \pracaq{d}: an antiquark (a bar over the quark letter).
\providecommand\pracaq[1]{\ensuremath{\overline{\mathrm{#1}}}}

\unnumberedchapter[\faDumbbell]{Practice set}

These 25 questions are written in the style of the HSC exam. They test all five chapters of
this book. They start easy and get harder.
\begin{itemize}
  \item \textbf{Section I:} 9 multiple-choice questions, 1 mark each. Choose A, B, C or D.
  \item \textbf{Section II:} 10 short-answer questions, 3 or 4 marks each.
  \item \textbf{Section III:} 6 extended-response questions, 5 to 8 marks each.
\end{itemize}
The total is 81 marks. Allow about 2 hours and 25 minutes. (In the HSC exam you have about
1.8 minutes for each mark.)

\textbf{For every question:}
\begin{itemize}
  \item Use the formula sheet and these data-sheet values: $e = \qty{1.602e-19}{C}$,
        $m_\mathrm{e} = \qty{9.109e-31}{kg}$, $m_\mathrm{p} = \qty{1.673e-27}{kg}$,
        $h = \qty{6.626e-34}{J.s}$, $c = \qty{3.00e8}{m.s^{-1}}$,
        $R = \qty{1.097e7}{m^{-1}}$, $b = \qty{2.898e-3}{m.K}$, $g = \qty{9.8}{m.s^{-2}}$,
        $\qty{1}{\amu} = \qty{931.5}{MeV}/c^2$ and $\qty{1}{eV} = \qty{1.602e-19}{J}$
        (so $\qty{1}{MeV} = \qty{1.602e-13}{J}$).
  \item Show your working. Give each answer with a unit and sensible significant figures.
\end{itemize}

The marks for each part of a question are shown on the right, like this:
\textcolor{hscmuted}{\small(2 marks)}. The tag next to each question number (for example
\textcolor{hscbrand}{Ch~2}) shows which chapter the question tests. If you get a question
wrong, revise that chapter. Full worked answers and marking criteria are at the back of the
book. Click a question number to jump to its answer.

% =====================================================================
\Needspace{15\baselineskip}
\section*{Section I: Multiple choice}

\Needspace{12\baselineskip}
\begin{practiceq}{1}{1}\label{pq:bigbang}
Less than a second after the Big Bang, the Universe was filled with very high-energy photons.
Which statement best describes how the first matter formed?
\choices[1]{Photons joined together and made hydrogen atoms straight away.}
  {Matter appeared from nothing, so energy was not conserved.}
  {Photons with enough energy turned into pairs of particles and antiparticles, as
   $E = mc^2$ allows.}
  {The first stars formed, and their fusion made the first protons and neutrons.}
\end{practiceq}

\Needspace{9\baselineskip}
\begin{practiceq}{1}{4}\label{pq:joliot}
In 1934 Ir\`ene Curie and Fr\'ed\'eric Joliot fired alpha particles at aluminium. They made
the first artificial (made by people) radioactive isotope:
\[
  \pracNuc{27}{13}{Al} + \pracNuc{4}{2}{He} \to \pracNuc{30}{15}{P} + \mathrm{X}
\]
What is X?
\choices{a proton}{a neutron}{an electron}{a gamma ray}
\end{practiceq}

\Needspace{8\baselineskip}
\begin{practiceq}{1}{5}\label{pq:dmeson}
The $\mathrm{D}^+$ particle is made of a charm quark and an anti-down quark
($\mathrm{c}\,\pracaq{d}$). What is its charge, and is it a fermion or a boson?
\choices[2]{charge $+1e$; a boson}{charge $+1e$; a fermion}
  {charge $+\tfrac13 e$; a boson}{charge $+\tfrac13 e$; a fermion}
\end{practiceq}

\Needspace{11\baselineskip}
\begin{practiceq}{1}{2}\label{pq:selector}
A beam of electrons passes straight between two charged metal plates. The plates make an
electric field. There is also a magnetic field. The two fields are at right angles to each
other and to the beam. The electrons are then replaced by protons with the same speed and
direction. The fields do not change. What happens to the protons?
\choices[1]{They bend towards the negative plate, because the electric force on them is
   reversed.}
  {They bend less than the electrons would, because they have a larger mass.}
  {They move in a circle, because the magnetic force on them is larger.}
  {They pass straight through, because both forces are reversed and still balance.}
\end{practiceq}

\Needspace{12\baselineskip}
\begin{practiceq}{1}{2}\label{pq:chadwick}
In 1932 Chadwick showed that the neutral radiation from beryllium was not gamma rays. Which
argument did he use to show this?
\choices[1]{Electric and magnetic fields did not bend the radiation, so it had no charge.}
  {The radiation was stopped by a thin sheet of paper, so it was made of heavy particles.}
  {To knock such fast protons out of paraffin wax, gamma photons would need more energy than
   the reaction could supply.}
  {Its speed was measured directly, and it was much less than the speed of light.}
\end{practiceq}

\Needspace{8\baselineskip}
\begin{practiceq}{1}{3}\label{pq:longest}
Electrons in hydrogen atoms make the four jumps below. Which jump gives out the photon with
the \emph{longest} wavelength?
\choices{$n = 2 \to n = 1$}{$n = 3 \to n = 2$}{$n = 4 \to n = 2$}{$n = 4 \to n = 3$}
\end{practiceq}

\Needspace{13\baselineskip}
\begin{practiceq}{1}{4}\label{pq:becurve}
The fission of uranium-235 and the fusion of hydrogen-2 both release energy. Which statement
explains why?
\choices[1]{In both reactions the products have a smaller binding energy per nucleon, so
   energy is released as the nuclei break apart.}
  {In both reactions the products have a larger binding energy per nucleon, so the products
   have less mass than the starting nuclei.}
  {In both reactions some nucleons disappear, and their mass becomes energy.}
  {Fission releases energy because the products have more mass. Fusion releases energy
   because the product has less mass.}
\end{practiceq}

\Needspace{11\baselineskip}
\begin{practiceq}{1}{5}\label{pq:cyclotronmc}
The ions in a cyclotron move much more slowly than light. Which statement about the
cyclotron is correct?
\choices[1]{The electric field in the gap speeds the ions up, and the magnetic field bends
   their path.}
  {The magnetic field does work on the ions, so they speed up on every half turn.}
  {The radius of the ions' path stays the same as they speed up.}
  {The frequency of the AC supply must be increased as the ions speed up.}
\end{practiceq}

\Needspace{12\baselineskip}
\begin{practiceq}{1}{1}\label{pq:magnitude}
Stars X and Y have the same surface temperature. The absolute magnitude of X is $-2.5$. The
absolute magnitude of Y is $+7.5$. Which statement is correct?
\choices[1]{Y is $10^4$ times more luminous than X, because its magnitude is larger.}
  {X is 10 times more luminous than Y, and the two stars have the same radius.}
  {X is $10^4$ times more luminous than Y, and X has a much larger radius than Y.}
  {X is $10^4$ times more luminous than Y, and the two stars have the same radius.}
\end{practiceq}

% =====================================================================
\section*{Section II: Short answer}

\Needspace{10\baselineskip}
\begin{practiceq}{3}{2}\label{pq:millikan}
An oil drop has a mass of \qty{4.2e-15}{kg}. It is held still between two horizontal metal
plates that are \qty{15.0}{mm} apart. The top plate is positive. The voltage between the
plates is \qty{964}{V}.
\begin{enumerate}[label=(\alph*)]
  \item Calculate the size of the charge on the drop. \pracmarks{2 marks}
  \item How many extra electrons does the drop carry? Justify your answer.
        \pracmarks{1 mark}
\end{enumerate}
\end{practiceq}

\Needspace{10\baselineskip}
\begin{practiceq}{3}{4}\label{pq:oxygen}
An oxygen-16 nucleus ($\pracNuc{16}{8}{O}$) has a mass of \qty{15.990523}{\amu}. A proton
has a mass of \qty{1.007276}{\amu}, and a neutron has a mass of \qty{1.008665}{\amu}.
\begin{enumerate}[label=(\alph*)]
  \item Calculate the mass defect of oxygen-16. \pracmarks{1 mark}
  \item Calculate its binding energy in MeV. \pracmarks{1 mark}
  \item Calculate its binding energy per nucleon. \pracmarks{1 mark}
\end{enumerate}
\end{practiceq}

\Needspace{10\baselineskip}
\begin{practiceq}{4}{1}\label{pq:wien}
The continuous spectrum of a star peaks at a wavelength of \qty{460}{nm}.
\begin{enumerate}[label=(\alph*)]
  \item Calculate the surface temperature of the star. \pracmarks{1 mark}
  \item State the star's spectral class and its colour. \pracmarks{1 mark}
  \item The star's spectrum also has dark absorption lines. Explain how these lines are
        produced, and how this spectrum is different from a black body spectrum.
        \pracmarks{2 marks}
\end{enumerate}
\end{practiceq}

\Needspace{10\baselineskip}
\begin{practiceq}{3}{3}\label{pq:debroglie}
Low-energy electrons are used to study the surface of a crystal, as in Davisson and Germer's
experiment. To give a clear diffraction pattern, the electrons need a de Broglie wavelength
of \qty{1.10e-10}{m}, close to the spacing of the atoms. The electrons start from rest.
\begin{enumerate}[label=(\alph*)]
  \item Calculate the speed the electrons need. \pracmarks{1 mark}
  \item Calculate the voltage needed to accelerate the electrons to this speed.
        \pracmarks{2 marks}
\end{enumerate}
\end{practiceq}

\Needspace{12\baselineskip}
\begin{practiceq}{4}{1}\label{pq:hubble}
The hydrogen-beta line has a wavelength of \qty{486.1}{nm} in the laboratory. In the spectrum
of a distant galaxy, this line is at \qty{497.8}{nm}. Take
$H_0 = \qty{68.0}{km.s^{-1}.Mpc^{-1}}$.
\begin{enumerate}[label=(\alph*)]
  \item Calculate the recession velocity of the galaxy. \pracmarks{2 marks}
  \item Calculate the distance to the galaxy in megaparsecs. \pracmarks{1 mark}
  \item Almost every distant galaxy shows a redshift. A student says that this means the
        Milky Way is at the centre of the Universe. Explain why this conclusion is wrong.
        \pracmarks{1 mark}
\end{enumerate}
\end{practiceq}

\Needspace{12\baselineskip}
\begin{practiceq}{3}{1}\label{pq:ppstep}
The second step of the proton--proton chain is
\[
  \pracNuc{2}{1}{H} + \pracNuc{1}{1}{H} \to \pracNuc{3}{2}{He} + \gamma
\]
The atomic masses are: hydrogen-2 \qty{2.014102}{\amu}, hydrogen-1 \qty{1.007825}{\amu}
and helium-3 \qty{3.016029}{\amu}.
\begin{enumerate}[label=(\alph*)]
  \item Calculate the energy released, in MeV. \pracmarks{2 marks}
  \item Assume that the gamma photon carries all of this energy. Calculate its wavelength.
        \pracmarks{1 mark}
\end{enumerate}
\end{practiceq}

\Needspace{13\baselineskip}
\begin{practiceq}{4}{2}\label{pq:thomson}
In a Thomson-type tube, a beam of electrons passes between two horizontal plates. The plates
are \qty{2.0}{cm} apart, with \qty{400}{V} between them. A magnetic field of
\qty{8.0e-4}{T} acts at right angles to the electric field and to the beam. The beam goes
straight.
\begin{enumerate}[label=(\alph*)]
  \item Calculate the speed of the electrons. \pracmarks{2 marks}
  \item The electric field is switched off. The same magnetic field bends the beam into a
        circular arc of radius \qty{17.8}{cm}. Calculate the charge-to-mass ratio of the
        electrons. \pracmarks{1 mark}
  \item Millikan measured the charge of the electron as \qty{1.602e-19}{C}. Use your answer
        to (b) to calculate the mass of the electron. \pracmarks{1 mark}
\end{enumerate}
\end{practiceq}

\Needspace{12\baselineskip}
\begin{practiceq}{4}{4}\label{pq:technetium}
Technetium-99m is used in medical scans. Its half-life is \qty{6.01}{h}. A patient is
injected with a dose that has an activity of \qty{6.0e8}{Bq}.
\begin{enumerate}[label=(\alph*)]
  \item Calculate the decay constant in \si{s^{-1}}. \pracmarks{1 mark}
  \item Calculate the number of technetium-99m nuclei in the dose when it is injected.
        \pracmarks{1 mark}
  \item Calculate the activity of the dose \qty{24.0}{h} later. \pracmarks{1 mark}
  \item Calculate how long it takes for the activity to fall to \qty{1.0}{\percent} of its
        starting value. Give your answer in hours. \pracmarks{1 mark}
\end{enumerate}
\end{practiceq}

\Needspace{16\baselineskip}
\begin{practiceq}{4}{4}\label{pq:fission}
One possible fission of uranium-235 is
\[
  \pracNuc{235}{92}{U} + \pracNuc{1}{0}{n} \to \pracNuc{141}{56}{Ba} + \pracNuc{92}{36}{Kr}
  + x\,\pracNuc{1}{0}{n}
\]
\begin{table}
  \caption*{Binding energy per nucleon of the three nuclei.}
  \centering
  \begin{tabular}{@{}lccc@{}}
    \toprule
    Nucleus & uranium-235 & barium-141 & krypton-92\\
    \midrule
    Binding energy per nucleon (\si{MeV}) & 7.59 & 8.33 & 8.51\\
    \bottomrule
  \end{tabular}
\end{table}
\begin{enumerate}[label=(\alph*)]
  \item Find $x$. \pracmarks{1 mark}
  \item Use the table to estimate the energy released by one fission. \pracmarks{2 marks}
  \item In a nuclear reactor, control rods keep this chain reaction steady. Explain how.
        \pracmarks{1 mark}
\end{enumerate}
\end{practiceq}

\Needspace{11\baselineskip}
\begin{practiceq}{4}{5}\label{pq:cyclotron}
A hospital cyclotron accelerates deuterons. A deuteron is a hydrogen-2 nucleus. Its charge is
$+e$ and its mass is \qty{3.344e-27}{kg}. The magnetic field is \qty{1.2}{T}. The deuterons
leave the cyclotron when the radius of their path is \qty{0.45}{m}. The voltage across the
gap between the dees is \qty{40}{kV}.
\begin{enumerate}[label=(\alph*)]
  \item Calculate the frequency of the AC supply. \pracmarks{1 mark}
  \item Calculate the kinetic energy of the deuterons when they leave, in MeV.
        \pracmarks{2 marks}
  \item Estimate how many turns a deuteron makes in the cyclotron. (Hint: in each turn it
        crosses the gap twice.) \pracmarks{1 mark}
\end{enumerate}
\end{practiceq}

% =====================================================================
\section*{Section III: Extended response}

\Needspace{10\baselineskip}
\begin{practiceq}{5}{2}\label{pq:goldfoil}
Thomson's model of the atom predicted that alpha particles fired at thin gold foil would
pass through with only very small deflections.
\begin{enumerate}[label=(\alph*)]
  \item Describe the results that Geiger and Marsden observed. Explain how Rutherford used
        these results to develop his nuclear model of the atom. \pracmarks{3 marks}
  \item Classical physics predicts that Rutherford's atom cannot be stable. Explain why.
        \pracmarks{2 marks}
\end{enumerate}
\end{practiceq}

\Needspace{12\baselineskip}
\begin{practiceq}{6}{3}\label{pq:balmer}
A student looks at the light from a hydrogen discharge tube through a spectrometer. Its
diffraction grating has 500 lines per millimetre. She turns the telescope to the
first-order cyan (blue-green) line. She reads the angle once, on one side of the
straight-through position only. Her reading is $\theta = \ang{14.1}$. The angle scale reads
to the nearest \ang{0.1}.
\begin{enumerate}[label=(\alph*)]
  \item Calculate the wavelength of the cyan line. \pracmarks{2 marks}
  \item Use the Rydberg equation to show that this line is produced when electrons in
        hydrogen drop from $n = 4$ to $n = 2$. \pracmarks{1 mark}
  \item Evaluate the reliability of her method, and suggest TWO improvements.
        \pracmarks{3 marks}
\end{enumerate}
\end{practiceq}

\begin{practiceq}{6}{4}\label{pq:decaygraph}
A student measures the count rate from a protactinium-234m source every \qty{30}{s}. The
background count rate is 30 counts per minute.
\begin{table}
  \caption*{Count rates from the protactinium-234m source.}
  \centering\small\setlength\tabcolsep{4pt}
  \begin{tabular}{@{}L{36mm}*{9}{C{8.5mm}}@{}}
    \toprule
    Time $t$ (\si{s}) & 0 & 30 & 60 & 90 & 120 & 150 & 180 & 210 & 240\\
    \midrule
    Measured count rate (counts per minute) & 1001 & 797 & 577 & 437 & 328 & 250 & 195 & 150 & 117\\
    \midrule
    Corrected count rate (counts per minute) & \pracblank & \pracblank & \pracblank &
      \pracblank & \pracblank & \pracblank & \pracblank & \pracblank & \pracblank\\
    \midrule
    $\ln(\text{corrected count rate})$ & \pracblank & \pracblank & \pracblank &
      \pracblank & \pracblank & \pracblank & \pracblank & \pracblank & \pracblank\\
    \bottomrule
  \end{tabular}
\end{table}
\begin{enumerate}[label=(\alph*)]
  \item The count rate decays in the same way as the number of nuclei,
        $N_\mathrm{t} = N_0 e^{-\lambda t}$. Show that a graph of
        $\ln N_\mathrm{t}$ against $t$ is a straight line with gradient $-\lambda$.
        (Hint: $\ln(ab) = \ln a + \ln b$.) \pracmarks{1 mark}
  \item Complete the table. (Corrected count rate = measured count rate $-$ background
        count rate.) \pracmarks{1 mark}
  \item On the grid below, plot $\ln(\text{corrected count rate})$ against $t$. Draw a line of
        best fit. \pracmarks{2 marks}
  \item Use your graph to calculate the half-life of protactinium-234m. \pracmarks{2 marks}
\end{enumerate}
\begin{figure}
  \centering
  \begin{tikzpicture}
    \begin{axis}[hsc, width=0.86\linewidth, height=0.6\linewidth,
        xmin=0, xmax=250, ymin=4.0, ymax=7.0,
        xlabel={Time $t$ (\si{s})}, ylabel={$\ln(\text{corrected count rate})$},
        xtick={0,50,...,250}, ytick={4.0,4.5,...,7.0},
        minor x tick num=4, minor y tick num=4,
        grid=both, minor grid style={hscrule!45, line width=0.2pt},
        yticklabel style={/pgf/number format/.cd, fixed, fixed zerofill, precision=1}]
      \addplot[draw=none] coordinates {(0,4)};
    \end{axis}
  \end{tikzpicture}
  \caption*{Grid for part (c).}
\end{figure}
\end{practiceq}

\begin{practiceq}{5}{3}\label{pq:matterwaves}
Bohr's model of the atom had limitations. Analyse how de Broglie's idea of matter waves
helped to explain one of these limitations, and how it led to Schr\"odinger's model.
\end{practiceq}

\begin{practiceq}{6}{5}\label{pq:microscope}
Particle accelerators are sometimes called ``the microscopes of particle physics''. Assess
this description. In your answer, refer to how accelerators work and to evidence that they
have produced for the Standard Model.
\end{practiceq}

\Needspace{28\baselineskip}
\begin{practiceq}{8}{1}\label{pq:hrdiagram}
The Hertzsprung--Russell (HR) diagram below shows four stars, P, Q, R and S, and the Sun. The
grey dots are other stars.
\begin{figure}
  \centering
  \begin{tikzpicture}
    \begin{axis}[hsc, width=0.84\linewidth, height=0.62\linewidth,
        xmode=log, ymode=log, x dir=reverse,
        xmin=2400, xmax=42000, ymin=1e-5, ymax=1e6,
        xtick={40000,20000,10000,5000,2500},
        xticklabels={\num{40000},\num{20000},\num{10000},5000,2500},
        ytick={1e-4,1e-2,1,1e2,1e4,1e6},
        xlabel={Surface temperature (\si{K})}, ylabel={Luminosity ($L_\odot$)}]
      \addplot[only marks, mark=*, mark size=1.3pt, color=hscmuted!55] coordinates {
        (3133,0.0164) (3434,0.0128) (3630,0.0506) (3819,0.0598) (4232,0.0655) (4283,0.0562)
        (4707,0.255) (5083,0.79) (5241,0.558) (5983,1.83) (6348,2.95) (6515,4.62)
        (7221,4.34) (7880,7.73) (8068,12) (9058,77.4) (9325,42.4) (10271,37.4) (10655,114)
        (12046,93.9) (12849,165) (13095,341) (14019,502) (15516,785) (15955,650)
        (17474,958) (19409,1870) (20326,4150) (22133,25000) (22629,12900) (25520,18600)
        (27432,26200) (28559,34900) (29945,49100)
        (3726,848) (3754,407) (5073,808) (4020,92.8) (4344,437) (3594,396) (4835,596)
        (3466,819) (3564,88.1) (4499,740)
        (3700,39800) (4300,79400) (9000,126000) (15000,200000) (6500,63100)
        (9747,0.000601) (9371,0.000236) (9430,0.000342) (7420,0.000143) (16043,0.00347)
        (6430,6.77e-05) (24530,0.0091) (8419,0.000181)};
      \addplot[only marks, mark=*, mark size=2.6pt, color=hscseriesA] coordinates {
        (3600,300) (20000,3000) (9000,0.003) (4200,0.15)};
      \addplot[only marks, mark=*, mark size=2.6pt, color=hscseriesB] coordinates {(5800,1)};
      \node[anchor=west, font=\small\bfseries, text=hscink] at (axis cs:3450,300) {P};
      \node[anchor=south, font=\small\bfseries, text=hscink, yshift=3pt] at (axis cs:20000,3000) {Q};
      \node[anchor=west, font=\small\bfseries, text=hscink, xshift=2pt] at (axis cs:9000,0.003) {R};
      \node[anchor=north, font=\small\bfseries, text=hscink, yshift=-3pt] at (axis cs:4200,0.15) {S};
      \node[anchor=east, font=\small, text=hscseriesB, xshift=-2pt] at (axis cs:5800,1) {Sun};
    \end{axis}
  \end{tikzpicture}
  \caption*{HR diagram. Both axes use log scales. Temperature increases to the left.}
\end{figure}
\begin{enumerate}[label=(\alph*)]
  \item Identify which star is a white dwarf and which is a red giant. Justify each choice
        using the diagram. \pracmarks{2 marks}
  \item Stars Q and S are both main sequence stars. Compare Q and S. Refer to their masses,
        how long they stay on the main sequence, and the main fusion process in their cores.
        \pracmarks{3 marks}
  \item The Sun will one day become a star like P. Describe the changes in its core, and
        the nuclear reactions that happen, from the end of its main sequence stage to the
        end of its life. \pracmarks{3 marks}
\end{enumerate}
\end{practiceq}
